%% kapa: exam=1 round=24 year=2013 period=2 form=A subject=accounting subject_ko=회계학 seq=19 number=19
%% source: Q-Net 감정평가사 기출문제 (cst00309) artlSeq=5073802 file=extracted/1008/5073802/2013년 제24회 감정평가사 1차시험 기출문제/2교시(A형).hwp
%% conversion: kordoc 4.24.51 -> tools/gen1.py (auto); review=proofread-sonnet; flags=none
\begin{question}{19}
(주)한국은 20×1년 1월 1일 만기 3년, 액면 ￦10,000의 전환사채를 액면발행하였다. 전환사채의 표시이자율은 연 7\%이고 이자는 매년 말에 지급한다. 전환조건은 다음과 같다.
\begin{framedtext}
○ 사채액면 ￦10당 1주의 보통주(액면가액 ￦5)로 전환 \par
○ 전환권이 행사되지 않은 부분에 대해서는 액면금액의 110\%를 일시 상환 \par
\end{framedtext}
발행시점에 전환권이 부여되지 않은 동일한 조건의 일반사채 시장이자율은 연 11\%이었다. 20×2년 1월 1일 사채 액면금액의 35\%가 전환되었을 경우, 전환권 행사가 20×2년 1월 1일 (주)한국의 재무상태표 상 자본총계에 미치는 영향은? (단, 이자율 11\%의 3년에 대한 단일금액 ￦1의 현가계수와 정상연금 ￦1의 현가계수는 각각 0.7312와 2.4437이며, 단수차이로 인한 오차가 있으면 가장 근사치를 선택한다.) \par
\begin{choices}
  \item ￦86 증가
  \item ￦1,750 증가
  \item ￦1,794 증가
  \item ￦1,880 증가
  \item ￦3,544 증가
\end{choices}
\end{question}
\answer{5}
